\documentclass[11pt,a4paper]{article}

\usepackage[utf8]{inputenc}
\usepackage[T1]{fontenc}
\usepackage{lmodern}
\usepackage{microtype}
\usepackage[margin=2.5cm]{geometry}
\usepackage{amsmath,amssymb,amsthm}
\usepackage{mathtools}
\usepackage{booktabs}
\usepackage{array}
\usepackage{graphicx}
\usepackage{enumitem}
\usepackage{hyperref}
\usepackage{xcolor}
\usepackage{tcolorbox}
\usepackage{fancyhdr}

\hypersetup{colorlinks=true, linkcolor=blue!50!black,
  citecolor=blue!50!black, urlcolor=blue!50!black}

\pagestyle{fancy}
\fancyhf{}
\fancyhead[L]{\small Icosian realization of $\mathbb{Q}(\sqrt5)$ Hilbert modular forms}
\fancyhead[R]{\small Pre-peer-review preprint}
\fancyfoot[C]{\thepage}

\newtheorem{theorem}{Theorem}
\newtheorem{proposition}[theorem]{Proposition}
\newtheorem{lemma}[theorem]{Lemma}
\newtheorem{corollary}[theorem]{Corollary}
\newtheorem{definition}[theorem]{Definition}
\newtheorem{conjecture}[theorem]{Conjecture}

\newcommand{\Q}{\mathbb{Q}}
\newcommand{\Z}{\mathbb{Z}}
\newcommand{\R}{\mathbb{R}}
\newcommand{\C}{\mathbb{C}}
\newcommand{\F}{\mathbb{F}}
\newcommand{\PP}{\mathbb{P}}
\newcommand{\Ord}{\mathcal{O}}
\newcommand{\fp}{\mathfrak{p}}
\newcommand{\fq}{\mathfrak{q}}
\newcommand{\Brandt}{B}
\newcommand{\OK}{\Ord_K}
\DeclareMathOperator{\Tr}{Tr}
\DeclareMathOperator{\Nm}{N}
\DeclareMathOperator{\disc}{disc}

\newtcolorbox{verifiedbox}[1][]{
  colback=green!3, colframe=green!45!black,
  fonttitle=\bfseries, title={Verified (computational, reproducible)}, #1}
\newtcolorbox{opengapbox}[1][]{
  colback=red!2, colframe=red!50!black,
  fonttitle=\bfseries, title={Open / not claimed}, #1}
\newtcolorbox{scopebox}[1][]{
  colback=gray!4, colframe=gray!55!black,
  fonttitle=\bfseries, title={Scope}, #1}

\title{\textbf{A native icosian realization of Hilbert modular forms
over $\Q(\sqrt5)$}\\[0.4em]
\large The $600$-cell unit group reproduces a family of cuspidal Hecke
systems --- parameter-free and out-of-sample}

\author{%
  Lee Smart\\[4pt]
  \textit{Institute of Vibrational Field Dynamics}\\[2pt]
  \texttt{contact@vibrationalfielddynamics.org}\\[2pt]
  \texttt{@vfd\_org}%
}

\date{2026-06-12 \quad (pre-peer-review preprint, v1.0.0-rc2)}

\begin{document}
\maketitle

\begin{abstract}
\noindent
We give an independent, parameter-free implementation of the standard
definite Brandt method for cuspidal Hilbert modular forms over
$K=\Q(\sqrt5)$, driven entirely by the geometry of the $600$-cell. The
binary icosahedral group $2I$ is the reduced-norm-one unit group of a
maximal order $\Ord$ (the \emph{icosian ring}) in the definite
quaternion algebra ramified at the two infinite places of $K$; its
lattice is $E_8$. For a split prime $\fp\subset\OK$ we build the Brandt
module of $\Ord$ at level $\fp$ purely from the action of
$A_5=2I/\{\pm1\}$ on $\PP^1(\OK/\fp)$, with no input from elliptic
curves, modular symbols, or $L$-function tables. The resulting Hecke
operators are integral and self-adjoint, and we verify against
independently point-counted data that this construction reproduces
\emph{a family} of forms, not a single example:
(i) the complete prime-level cuspidal-dimension sequence over $K$ for
norm ${<}100$ ($7/7$ agreements with Demb\'el\'e, including the two
non-obvious cases --- the absent form at norm $59$ and the jump to a
genus-$2$ form at norm $61$);
(ii) the full Hecke eigenvalue sets of \emph{two} rational eigenforms
(levels $31$ and $41$), matched out-of-sample against brute-force point
counts of the associated elliptic curves at every prime tested, with no
fitting;
(iii) the genus-$2$ form at level $61$, whose irrational eigenvalues in
$\Z[\varphi]$ (real multiplication by $K$) the geometry reconstructs,
matching the published values exactly at the inert primes.
We close with an explicitly heuristic discussion of how this finite,
self-adjoint, coefficient-side structure relates to the Riemann
Hypothesis, emphasizing the gap rather than bridging it: RH is the
\emph{global, analytic} positivity equivalent to the Weil and Li
criteria, of which a finite Brandt computation supplies nothing.
\textbf{This paper proves no Millennium problem and claims no progress
on one.} Its contribution is a clean, reproducible geometric realization
of known arithmetic, and an honest delimitation of the gap that remains.
\end{abstract}

\begin{scopebox}
This is a \emph{reproduction-and-map} paper. The arithmetic it recovers
is due to Demb\'el\'e \cite{Dembele2005} and is an instance of the
Jacquet--Langlands correspondence \cite{JacquetLanglands1970,
Eichler1973}; our contribution is the independent, parameter-free
\emph{geometric} route (native icosian arithmetic, no computer-algebra
modular package), the explicit no-fit out-of-sample protocol, and a
precise statement of what the verified structure does \emph{not} reach
(\S\ref{sec:rh}). The arithmetic targets are Demb\'el\'e's known
examples; the novelty is independent implementation and provenance,
not new forms. We claim \emph{validation and exposition}, not new
number theory, and \emph{no} progress on RH itself.
\end{scopebox}

\section{Introduction}\label{sec:intro}

The $600$-cell is the regular convex $4$-polytope whose $120$ vertices
are the unit icosians: the binary icosahedral group $2I\subset
\mathbb{H}$ of order $120$. Viewed as a lattice in $\R^4\cong\mathbb{H}$,
the icosians are a copy of $E_8$ over $\Z$, and a copy of the Hurwitz-type
maximal order $\Ord$ in the definite quaternion algebra
$B=\left(\tfrac{-1,-1}{K}\right)$ over $K=\Q(\sqrt5)$
\cite{ConwaySloane1999,Voight2021}. This is the smallest arithmetic
setting in which the golden ratio $\varphi=(1+\sqrt5)/2$ and the
icosahedral symmetry that organizes the $600$-cell become number
theory.

By the Jacquet--Langlands correspondence and the Eichler basis problem
\cite{JacquetLanglands1970,Eichler1973,Vigneras1980}, the Brandt module
of $\Ord$ at a prime level $\fp$ computes the Hilbert cusp forms of
parallel weight $2$ and level $\fp$ over $K$. Demb\'el\'e
\cite{Dembele2005} used exactly this to produce the first explicit
tables of such forms; the smallest level carrying a cusp form has norm
$31$, and the associated rational eigenform corresponds to an elliptic
curve over $K$ \cite{BoberEtAl2012}.

Our question is narrow and testable: \emph{does the $600$-cell geometry,
used on its own terms, reproduce this arithmetic?} We answer yes, and
quantitatively. The validation is anti-circular by construction: every
target we compare against is generated independently --- by brute-force
point counting of an elliptic curve, or from the published literature
--- and is used \emph{only in comparison}, never in building the Brandt
operators.

\paragraph{Results.}
\begin{enumerate}[leftmargin=2em]
\item \textbf{One form, completely (\S\ref{sec:level31}).} At level
$\fp_{31}$ the geometry gives a $2$-dimensional Brandt module, Eichler
mass $8/15$, a $1$-dimensional cuspidal eigenform, and Hecke
eigenvalues matching the point-counted curve $31.1\text{-}a1$ at all
$44$ good prime ideals of norm ${\le}200$, \emph{individually}: the
splitting choice $\varphi\mapsto13$ places the engine's level at
$\sigma(\fp_{31})$, so the geometric system is the Galois conjugate of
the curve's; after fixing that single global involution once, every
labeled ideal matches with no per-prime freedom, and the Brandt value
at the engine's own level ideal recovers the Steinberg datum $-1$ used
by the Atkin--Lehner check.
\item \textbf{A family, by dimension (\S\ref{sec:dims}).} The geometric
orbit count reproduces the entire prime-level cuspidal-dimension
sequence over $K$ for norm ${<}100$: $7/7$ agreement with
\cite{Dembele2005}, including the absence of a form at norm $59$ and
the jump to cuspidal dimension $2$ at norm $61$.
\item \textbf{A family, by eigenvalue (\S\ref{sec:level41}).} A second
rational eigenform, at level $\fp_{41}$, is built by the same procedure
and matched out-of-sample (full eigenvalue sets, $9/9$ primes) against
the point-counted curve $E_{41}:y^2+\varphi y=x^3-\varphi x^2$.
\item \textbf{The non-rational form (\S\ref{sec:level61}).} At level
$\fp_{61}$ the geometry yields cuspidal dimension $2$ --- the
genus-$2$ abelian surface with real multiplication by $K$ of
\cite[Rem.~4.4]{Dembele2005}. From the $3\times3$ Brandt matrix we
recover, at each of the six tested primes (norms $4$--$49$), an
eigenvalue whose field-theoretic discriminant is
$5\cdot(\text{square})$, i.e.\ an irrational element of $\Z[\varphi]$;
at the three inert primes the values match the published ones exactly.
\item \textbf{The certificate algebra, and a precise non-claim
(\S\ref{sec:cert}--\S\ref{sec:rh}).} The Brandt operators are
self-adjoint for the orbit-size measure --- an exactly checkable
certificate in a small algebra we make precise. A closing section
states exactly which analytic statements (Weil/Li positivity, i.e.\
GRH-type questions for the assembled $L$-functions) the finite
computation does \emph{not} reach, and why a finite coefficient
computation cannot reach them.
\end{enumerate}

All computations are pure Python, reproduce in seconds to minutes, and
are included with the paper; \S\ref{sec:repro} gives the manifest.

\section{The icosian order and its Brandt module}\label{sec:setup}

Let $K=\Q(\sqrt5)$, $\varphi=(1+\sqrt5)/2$, $\OK=\Z[\varphi]$, with
$\Nm(a+b\varphi)=a^2+ab-b^2$ and nontrivial automorphism
$\sigma:\varphi\mapsto1-\varphi$. Let $B=\left(\tfrac{-1,-1}{K}\right)$
with standard basis $1,i,j,k$. The \emph{icosian ring} $\Ord$ is the
maximal order whose group of reduced-norm-one units is the binary
icosahedral group, $\Ord^1=2I$, $|2I|=120$; as a $\Z$-lattice $\Ord$ is
$E_8$ \cite{ConwaySloane1999,Voight2021}. $B$ is ramified exactly at the
two infinite places of $K$, hence definite, and (its maximal orders
form) a single ideal class.

Throughout we take the level $\fp\subset\OK$ to be a prime that
\emph{splits} in $K$, so that $\Nm(\fp)=q$ is a rational prime and the
residue field is $\OK/\fp\cong\F_q$. (The levels $31,41,61$ used below
are all split; the inert and ramified cases, with residue field
$\F_{q^2}$ or $\F_q$, are not needed.) For odd $q$ fix a local splitting
\[
  \iota_\fp:\ \Ord\otimes_{\OK}(\OK)_\fp\ \xrightarrow{\ \sim\ }\
  M_2(\OK/\fp)=M_2(\F_q),
\]
constructed explicitly by solving $i^2=j^2=-1$, $ij=-ji$ over $\F_q$
(possible since $-1$ is a sum of two squares mod $q$). The unit group
$2I$ acts through $\iota_\fp$ on the $q+1$ points of $\PP^1(\F_q)$;
since $-1\in2I$ acts trivially, the action factors through
$A_5=2I/\{\pm1\}$. Because $\Ord$ has class number one, the right-ideal
classes of the Eichler order $\Ord_0(\fp)$ identify with
$\Ord^1\backslash\PP^1(\OK/\fp)$, i.e.\ with the $A_5$-orbits on
$\PP^1(\F_q)$ (the definite/Brandt method;
\cite[Ch.~41]{Voight2021}, \cite{Vigneras1980}). Write $h$ for the
number of orbits and $\mu=(\mu_1,\dots,\mu_h)$ for their sizes, so
$\sum_i\mu_i=q+1$.

\paragraph{Brandt operators.}
For a totally positive $\varpi\in\OK$ with $\Nm(\varpi)=\Nm(\fq)$,
$\fq\nmid\fp$, enumerate the icosians of reduced norm $\varpi$ (a
short-vector enumeration in $E_8$). Their number is
$120\,r(\varpi)$ where $r(\varpi)=\Nm(\fq)+1$ is the relevant Brandt /
theta coefficient for the icosian order at the primes we use; we verify
this count at construction time as a gate. Push each element through
$\iota_\fp$ to a matrix in $\mathrm{GL}_2(\F_q)$, and let it permute
$\PP^1(\F_q)$. Accumulating the induced map on orbits gives the Brandt
matrix $\Brandt(\fq)\in M_h(\Q)$,
\begin{equation}
  \Brandt(\fq)_{IJ}=\frac{1}{120\,\mu_I}
  \#\{\,(x,\gamma): x\in\text{orbit }I,\ \Nm(\gamma)=\varpi,\
  x\cdot\iota_\fp(\gamma)\in\text{orbit }J\,\}.
  \label{eq:brandt}
\end{equation}
These are the Hecke operators $T_\fq$ on the Brandt module. Three
intrinsic invariants are checked at construction time, before any
arithmetic target is consulted:
\begin{itemize}[leftmargin=1.5em]
\item \emph{Eisenstein row sums:} $\sum_J\Brandt(\fq)_{IJ}=\Nm(\fq)+1$
(the all-ones vector is the Eisenstein eigenvector);
\item \emph{integrality:} $\Brandt(\fq)\in M_h(\Z)$;
\item \emph{self-adjointness:} $\mu_I\Brandt(\fq)_{IJ}=
\mu_J\Brandt(\fq)_{JI}$, i.e.\ $\Brandt(\fq)$ is symmetric for the inner
product weighted by the orbit sizes $\mu$.
\end{itemize}
The cuspidal subspace is the $\mu$-orthogonal complement of the
Eisenstein vector; on it the eigenvalues of $\Brandt(\fq)$ are the
Hecke eigenvalues $a_\fq$ of the cusp forms of level $\fp$.

\begin{scopebox}
Equation \eqref{eq:brandt} and the bijection ``classes $=$ $A_5$-orbits''
are standard for class number one \cite{Vigneras1980,Voight2021}. The
point of the paper is that we instantiate them with nothing but icosian
short-vector arithmetic and the $A_5$-action, and then test the output
against independent data.
\end{scopebox}

\section{Level 31: one form, completely}\label{sec:level31}

At $\fp_{31}=(5\varphi-2)$, $\Nm(\fp_{31})=31$ (the smallest norm
carrying a cusp form \cite{Dembele2005}), the $A_5$-action on
$\PP^1(\F_{31})$ has two orbits, $\mu=(12,20)$, so $h=2$ and the
cuspidal dimension is $1$. The point stabilizers have orders
$60/12=5$ and $60/20=3$, giving Eichler mass
$\tfrac15+\tfrac13=\tfrac{8}{15}=\tfrac{1}{60}(31+1)$, in agreement with
the mass formula (with $|\zeta_K(-1)|/2=1/60$ for $K$). The cuspidal
eigenvector is $(3,-5)$.

\paragraph{Out-of-sample match.}
Independently of the geometry we form the target by brute-force point
counting of the elliptic curve $31.1\text{-}a1/K$,
$[a_1,\dots,a_6]=[1,\varphi+1,\varphi,\varphi,0]$, setting
$a_\fq=\Nm(\fq)+1-\#E(\OK/\fq)$. No $L$-function table is read.

A labeling convention must be fixed before a \emph{per-ideal} claim can
even be stated. Our level splitting sends $\varphi\mapsto13\pmod{31}$,
whose kernel is $\sigma(\fp_{31})$ --- the Galois conjugate of the
curve's conductor $\fp_{31}=(5\varphi-2)$ --- so the geometric Hecke
system is the $\sigma$-conjugate of the curve's: the Brandt eigenvalue
at a generator of $\sigma(\fq)$ equals $a_\fq$. With that single global
involution fixed once, the Brandt eigenvalues agree with $a_\fq$ at all
$44$ good prime ideals of norm $\le200$ \emph{individually} --- each
labeled ideal checked on its own, with no per-prime freedom anywhere
(\texttt{data/level31\_per\_ideal.json}; the verification script fails
if any single ideal needs its own adjustment). The engine's level ideal
is excluded from the good-prime comparison, and its Brandt value $-1$
is exactly the Steinberg datum consumed by the Atkin--Lehner sign check
below.

\begin{verifiedbox}
\textbf{Split-prime completeness.} For a split rational prime
$q=\fq\fq'$ the two Hecke eigenvalues $a_\fq,a_{\fq'}$ are the values at
the totally positive generator $\varpi$ and at its conjugate
$\sigma(\varpi)$. The geometry reproduces the \emph{full set}
$\{a_\fq,a_{\fq'}\}$ --- not one representative up to the global Galois
involution. For $q=11,\dots,89$:
\begin{center}\resizebox{\linewidth}{!}{$
\begin{array}{c|ccccccccc}
q & 11 & 19 & 29 & 41 & 59 & 61 & 71 & 79 & 89\\\hline
\{a_\fq,a_{\fq'}\} & \{\mbox{-}4,4\} & \{\mbox{-}4,4\} & \{\mbox{-}2\}
& \{\mbox{-}6\} & \{\mbox{-}4,12\} & \{\mbox{-}2,6\} & \{\mbox{-}8,0\}
& \{0,16\} & \{\mbox{-}6,10\}
\end{array}
$}\end{center}
matching the point-counted target exactly (singletons at $q=29,41$
occur where $a_\fq=a_{\fq'}$).
\end{verifiedbox}

\paragraph{Atkin--Lehner sign.}
\sloppy At $\fp_{31}$ the curve $31.1\text{-}a1$ has nonsplit multiplicative
reduction ($\#E_{\mathrm{ns}}(\F_{31})=q+1$, i.e.\ $a_{\fp_{31}}=-1$).
With the standard newform convention $W_\fp=-a_\fp$ at a prime of
multiplicative reduction exactly dividing the level, the Atkin--Lehner
(Fricke) eigenvalue is $W_{31}=-a_{\fp_{31}}=+1$, consistent with the
Brandt data.

\section{A family, by dimension}\label{sec:dims}

The orbit count of \S\ref{sec:setup} gives the Brandt dimension $h$, and
hence the cuspidal dimension $h-1$ at every prime level with no further
computation. (The Eisenstein subspace at prime level is $1$-dimensional
because $K=\Q(\sqrt5)$ has narrow class number one, so there is a single
Eisenstein class, the all-ones vector.) Table~\ref{tab:dims} compares
the geometric output with Demb\'el\'e's table \cite{Dembele2005} at
every split prime level in his displayed range ($31\le\Nm\le89$); the
smaller split prime norms ($11,19,29$) all give cuspidal dimension $0$,
consistent with $31$ being the smallest level carrying a cusp form.

\begin{table}[ht]
\centering
\begin{tabular}{r|ccccccc}
\toprule
$\Nm(\fp)$ & 31 & 41 & 59 & 61 & 71 & 79 & 89\\
\midrule
geometry (cusp.\ dim) & 1 & 1 & 0 & 2 & 1 & 1 & 1\\
Demb\'el\'e \cite{Dembele2005} & 1 & 1 & 0 & 2 & 1 & 1 & 1\\
type & ell. & ell. & --- & gen.~2 & ell. & ell. & ell.\\
\bottomrule
\end{tabular}
\caption{Prime-level cuspidal dimensions over $\Q(\sqrt5)$, geometry vs
Demb\'el\'e. $7/7$. The two informative cases are $59$ (no
weight-$(2,2)$ cusp form, only Eisenstein \cite[Ex.~3.6]{Dembele2005})
and $61$ (the first form with coefficients outside $\Q$, a genus-$2$
surface \cite[Rem.~4.4]{Dembele2005}).}
\label{tab:dims}
\end{table}

The point of Table~\ref{tab:dims} is that the orbit count is not tuned
to one example: with no free parameters it reproduces a non-obvious
\emph{sequence}, including a vanishing ($59$) and a jump to a
non-rational form ($61$). The same computation continues to larger
split levels at no extra cost; we do not tabulate those here, as we have
not cross-checked them against an independent source.

\section{A family, by eigenvalue: level 41}\label{sec:level41}

We repeat \S\ref{sec:level31} verbatim at $\fp_{41}$ ($\Nm=41$, split):
the $A_5$-action on $\PP^1(\F_{41})$ has two orbits, $\mu=(12,30)$,
$h=2$, cuspidal dimension $1$. The cuspidal eigenvalues are integral,
self-adjoint, and satisfy the Ramanujan bound
$|a_\fq|\le2\sqrt{\Nm(\fq)}$ at every prime tested.

For the out-of-sample target we point-count Demb\'el\'e's curve
$E_{41}:y^2+\varphi y=x^3-\varphi x^2$
(\cite[Tab.~4]{Dembele2005}, $[a_1,\dots,a_6]=[0,-\varphi,\varphi,0,0]$).

\begin{verifiedbox}
\textbf{Second form, out-of-sample, no fit.} Geometric eigenvalue sets
$\{a_\fq,a_{\fq'}\}$ vs point-counted $E_{41}$:
\begin{center}\resizebox{\linewidth}{!}{$
\begin{array}{c|ccccccccc}
q & 11 & 19 & 29 & 31 & 59 & 61 & 71 & 79 & 89\\\hline
\text{geom} & \{\mbox{-}2,5\} & \{\mbox{-}1,6\} & \{2,9\}
& \{\mbox{-}10,4\} & \{\mbox{-}3,4\} & \{\mbox{-}8,6\}
& \{\mbox{-}12,9\} & \{\mbox{-}11,\mbox{-}4\} & \{\mbox{-}8,\mbox{-}1\}\\
E_{41} & \{\mbox{-}2,5\} & \{\mbox{-}1,6\} & \{2,9\}
& \{\mbox{-}10,4\} & \{\mbox{-}3,4\} & \{\mbox{-}8,6\}
& \{\mbox{-}12,9\} & \{\mbox{-}11,\mbox{-}4\} & \{\mbox{-}8,\mbox{-}1\}
\end{array}
$}\end{center}
$9/9$. The same construction, re-pointed to a different level, recovers
a different eigenform eigenvalue-for-eigenvalue.
\end{verifiedbox}

\section{The non-rational form: level 61}\label{sec:level61}

At $\fp_{61}$ ($\Nm=61$, split) the $A_5$-action on $\PP^1(\F_{61})$ has
\emph{three} orbits, $\mu=(12,20,30)$, $h=3$, so the cuspidal dimension
is $2$. By \cite[Rem.~4.4]{Dembele2005}, $61$ is the smallest level
whose newform has coefficients outside $\Q$: the associated Shimura
curve has genus $2$, and its Jacobian is a modular abelian surface with
real multiplication by $\Q(\sqrt5)$. The $h=2$ shortcut of
\S\ref{sec:level31} does not apply; instead we read the eigenvalue off
the $3\times3$ matrix.

The all-ones vector is the Eisenstein eigenvector (eigenvalue
$\Nm(\fq)+1$), so the characteristic polynomial of $\Brandt(\fq)$
factors over $\Q$ as $(x-(\Nm(\fq)+1))\,g_\fq(x)$, where $g_\fq(x)=
x^2-t_\fq x+n_\fq$ is the characteristic polynomial of the $2$-dimensional
cuspidal block. Comparing traces and determinants,
\begin{equation}
  t_\fq=\operatorname{tr}\Brandt(\fq)-(\Nm(\fq)+1),
  \qquad
  n_\fq=\frac{\det\Brandt(\fq)}{\Nm(\fq)+1}\ \in\Z,
  \label{eq:rm}
\end{equation}
\emph{a priori} two rational integers. Whether the two cuspidal
eigenvalues are rational (two forms) or an irrational conjugate pair
(one form with real multiplication) is then decided by the
discriminant $t_\fq^2-4n_\fq$ --- which we compute, rather than assume.

\begin{verifiedbox}
\textbf{Real multiplication, reconstructed.} From \eqref{eq:rm} the
discriminant $t_\fq^2-4n_\fq$ equals $5\cdot(\text{square})$ at every
prime $\Nm(\fq)\in\{4,9,11,19,29,49\}$. Hence $g_\fq$ is irreducible
over $\Q$ with roots $a_\fq,\sigma(a_\fq)\in\Z[\varphi]$ a Galois-conjugate
pair: a \emph{single} eigenform with real multiplication by $K$, not two
rational forms. At the inert primes (one prime per rational prime, hence
unambiguous) the eigenvalue matches \cite{Dembele2005} exactly:
\[
  a_4=2\varphi-2,\qquad a_9=-\varphi-2,\qquad a_{49}=-4\varphi+2,
\]
i.e.\ $(t_\fq,n_\fq)=(-2,-4),(-5,5),(0,-20)$ respectively. The
Ramanujan bound holds on both real embeddings throughout.
\end{verifiedbox}

The real multiplication is not merely consistent with the data; it
follows from it:

\begin{proposition}\label{prop:rm}
If $g_\fq$ has non-square discriminant for even a single $\fq$, then the
$2$-dimensional cuspidal Brandt module at level $\fp_{61}$ is irreducible
over $\Q$, the (commutative) Hecke algebra acts on it through the
quadratic field $\Q(\sqrt{t_\fq^2-4n_\fq})$, and every Hecke eigenvalue
lies in that field. In particular, $t_4^2-4n_4=20=5\cdot2^2$ gives
$\Q(\sqrt5)$, so the level-$61$ eigenform has real multiplication by
$\Q(\sqrt5)$ and all its eigenvalues lie in $\Z[\varphi]$.
\end{proposition}
\begin{proof}
On the $2$-dimensional cuspidal space, $t_\fq^2-4n_\fq$ non-square makes
$g_\fq$ irreducible over $\Q$, so $T_\fq$ acts irreducibly. Any operator
commuting with such an irreducible quadratic operator lies in
$\Q[T_\fq]\cong\Q(\sqrt{t_\fq^2-4n_\fq})$, a field; the Hecke operators
all commute and preserve the space, so they lie in $\Q[T_\fq]$, and every
cuspidal Hecke eigenvalue lies in $E:=\Q(\sqrt{t_\fq^2-4n_\fq})$. The
value $t_4^2-4n_4=20$ fixes $E=\Q(\sqrt5)$; integrality of the Brandt
matrices puts the eigenvalues in $\OK=\Z[\varphi]$.
\end{proof}

So the geometry does not merely \emph{count} the dimension-$2$ space at
$61$, nor merely match six published values: it \emph{proves}, from a
single discriminant, that the form has real multiplication by
$\Q(\sqrt5)$, and reconstructs its irrational Hecke eigenvalues.

\section{The closure certificate algebra}\label{sec:cert}

The recurring structural fact above is that each Hecke operator carries
a positive invariant form. We isolate this as a small certificate
algebra, both because it is the cleanest statement of what the geometry
guarantees and because it is the bridge to \S\ref{sec:rh}.

\begin{definition}[Positive symmetrizer]
A \emph{positive symmetrizer} for $T\in M_n(\R)$ is a symmetric positive
definite $G$ with $GT=T^\top G$. Then $T$ is self-adjoint for the inner
product $\langle u,v\rangle_G=u^\top G v$, so it has real spectrum and a
$G$-orthonormal eigenbasis.
\end{definition}

\noindent
(One can place this within the usual operator taxonomy --- dissipative
$T^\top GT\prec G$, isometric $T^\top GT=G$ --- which we do not need
here: the Brandt operators are not contractions, having Eisenstein
eigenvalue $\Nm(\fq)+1>1$.)

\begin{proposition}\label{prop:selfadj}
For every $\fq\nmid\fp$ the Brandt operator $\Brandt(\fq)$ admits the
positive symmetrizer $G=\operatorname{diag}(\mu)$, hence has real
spectrum and a $\mu$-orthonormal eigenbasis.
\end{proposition}
\begin{proof}
The raw incidence count $C_{IJ}=\#\{(x,\gamma):x\in\text{orbit }I,
\Nm(\gamma)=\varpi, x\iota_\fp(\gamma)\in\text{orbit }J\}$ satisfies
$C_{IJ}=C_{JI}$ via the involution $(x,\gamma)\mapsto(x\gamma,
\bar\gamma)$ (using $\gamma\bar\gamma=\Nm(\gamma)$ acting as a scalar
on $\PP^1$). Since $\Brandt(\fq)_{IJ}=C_{IJ}/(120\,\mu_I)$ by
\eqref{eq:brandt}, this gives $\mu_I\Brandt(\fq)_{IJ}=
\mu_J\Brandt(\fq)_{JI}$, i.e.\ $G\Brandt(\fq)=\Brandt(\fq)^\top G$.
\end{proof}

\noindent
Real spectrum is thus elementary. The sharper \emph{Ramanujan bound}
$|a_\fq|\le2\sqrt{\Nm(\fq)}$ that we observe at every tested prime is
\emph{not} elementary: it is the Ramanujan--Petersson bound for these
Hilbert cusp forms, transported through Jacquet--Langlands and proved by
Eichler--Shimura/Deligne \cite{JacquetLanglands1970,Eichler1973}; in
this paper we record it as a property satisfied by the computed
eigenvalues, not as something we prove. The Brandt operators are, in
short, the \emph{positive, finite, real-spectrum} fragment of the
arithmetic.

\section{Scope: what the finite computation does not reach}\label{sec:rh}

Nothing in \S\ref{sec:level31}--\S\ref{sec:cert} depends on this
section, and it contains no new theorem. We include it only to state
precisely what the verified structure does \emph{not} reach.

\medskip\noindent\textbf{The rigorous statement.}
For an $L$-function, RH is equivalent to genuine positivity criteria:
the Weil explicit-formula positivity \cite{Weil1952} and Li's criterion
$\lambda_n\ge0$ for all $n$ \cite{Li1997,BombieriLagarias1999}. These
are theorems (equivalences); they concern the \emph{global, analytic}
$L$-function and an \emph{infinite} family of positivity conditions.

\medskip\noindent\textbf{What the geometry delivers, and what it does not.}
\begin{verifiedbox}
The icosian construction delivers the \emph{coefficient side}: the Hecke
eigenvalues $a_\fq$, equivalently the local factors of the automorphic
$L$-functions, with real spectrum (Prop.~\ref{prop:selfadj}) and the
observed Ramanujan bound. This is finite, per-prime, local data.
\end{verifiedbox}

\begin{opengapbox}
RH / GRH is the global, analytic statement that the assembled
$L$-function has no off-line zero --- the infinite positivity of
Weil/Li. A finite Brandt computation produces eigenvalue
\emph{coefficients}, not zeros, and supplies none of this positivity.
Programs that would close the gap --- a self-adjoint operator with the
zeros as spectrum (Hilbert--P\'olya), or Connes's arithmetic-site trace
formula \cite{Connes1999} --- remain open; we make no contribution to
them. Mirror symmetry (the functional equation) alone does not force
zeros onto the line: the Davenport--Heilbronn phenomenon
\cite{Davenport1936} exhibits Dirichlet series with the functional
equation yet off-line zeros.
\end{opengapbox}


\section{Reproducibility}\label{sec:repro}

All results are produced by pure-Python scripts in the accompanying
bundle (\texttt{route\_b/}); none requires a computer-algebra modular
package. The anti-circularity discipline is enforced in code: targets
are computed by \texttt{point\_count\_target.py} (brute-force
$\#E(\OK/\fq)$) and read only at comparison time.
\begin{center}
\small
\begin{tabular}{ll}
\toprule
script & result \\
\midrule
\texttt{ideal\_classes.py} & level-$31$ orbits, mass $8/15$, $h=2$ \\
\texttt{brandt\_matrices.py} & Brandt/Hecke operators \eqref{eq:brandt}, invariants \\
\texttt{brandt\_level.py} & the engine at an arbitrary split prime level \\
\texttt{hecke\_compare.py} & level $31$, original out-of-sample protocol \\
\texttt{level31\_full\_ideals.py} & level $31$, all $44$ ideals $\Nm\le200$
  $\to$ \texttt{data/level31\_per\_ideal.json} \\
\texttt{second\_form\_level41.py} & level $41$, $9/9$ vs $E_{41}$
  $\to$ \texttt{data/level41\_results.json} \\
\texttt{genus2\_form\_level61.py} & level $61$ genus-$2$, RM eigenvalues
  $\to$ \texttt{data/level61\_results.json} \\
\texttt{multilevel\_dimensions.py} & dimension sequence (Table~\ref{tab:dims})
  $\to$ \texttt{data/dimension\_sequence.json} \\
\texttt{no\_fit\_guard.py} & anti-circularity check over the full geometry path \\
\bottomrule
\end{tabular}
\end{center}

\section{Conclusion}

We have shown, by an independent geometric route, that the $600$-cell /
icosian order reproduces a \emph{family} of Hilbert modular forms over
$\Q(\sqrt5)$: a non-obvious dimension sequence ($7/7$), the full Hecke
systems of two rational eigenforms (out-of-sample, no fit), and the
irrational eigenvalues of the genus-$2$ form at level $61$. We have
recorded the self-adjointness of the Hecke operators as an exactly
checkable certificate, and used it to state exactly what separates this
verified, finite, coefficient-side structure from GRH-type questions:
the global analytic positivity of the assembled $L$-functions, which a
finite coefficient computation cannot supply. The value of the work is
a clean, reproducible realization of known arithmetic from polytope
geometry, and an honest account of the gap that remains.

\bibliographystyle{plain}
\bibliography{references}

\end{document}
